\bookchapter{פרק 40: הנדסת דרישות — מהפיכת רעיונות להוראות שה-AI יכול לבצע}

\chapterintro{אם הנדסת הקשר פותרת את השאלה “לגרום ל-AI להבין את הרקע”, הרי שהנדסת דרישות פותרת את השאלה “לגרום ל-AI לדעת מה בדיוק צריך לעשות הפעם”. זהו השלב שבו ניכר יותר מכל הפער בין אנשים — עם אותם כלים ואותו מודל, יש מי שמוציא פרויקטים מרהיבים ויש מי שמוציא רק צעצועים בינוניים. ההבדל בדרך כלל אינו ביכולת הטכנית, אלא ביכולת לנסח דרישות.}

\booksection{א. מרעיון מעורפל לדרישה מובנית}

כל פרויקט מתחיל ברעיון מעורפל — “רוצה לבנות אפליקציה לניהול חשבונות”, “רוצה אתר אישי”, “רוצה כלי לרשימת מטלות”. הרעיונות עצמם בסדר גמור, אבל אם תזרקו אותם ישירות אל ה-AI, איכות התוצאה תיקבע אך ורק על פי המזל. מטרת הנדסת הדרישות היא להפוך את הרעיון המעורפל הזה לדרישה מובנית שה-AI יכול לבצע במדויק.

החדשות הטובות: דרישה מובנית לא מחייבת מסמך PRD בן עשרות עמודים. היא מחייבת תשובה לחמש שאלות, שנקראות “חמשת מרכיבי הדרישה”:

\begin{center}
\begin{tabularx}{\linewidth}{|X|X|X|}
\hline
\textbf{מרכיב} & \textbf{שאלה} & \textbf{דוגמה: אפליקציית ניהול חשבונות} \tabularnewline
\hline
Who & למי זה מיועד? & שימוש אישי, לתיעוד הכנסות והוצאות יומיומיות \tabularnewline
What & מה עושים? & אפליקציה שמתעדת כל הכנסה והוצאה, מסכמת לפי חודש ומציגה ויזואליזציה \tabularnewline
Why & למה עושים את זה? & רוצה להבין את מבנה ההוצאות שלי ולזהות הוצאות מיותרות \tabularnewline
How & בגדול, איך עושים את זה? & פרונט-אנד ב-React, נתונים נשמרים מקומית בדפדפן, ובינתיים בלי בק-אנד והתחברות \tabularnewline
Constraint & מה המגבלות? & פיתוח על ידי אדם אחד, להשלים בתוך שבועיים, בלי פריסה לאוויר \tabularnewline
\hline
\end{tabularx}
\end{center}

\begin{center}\small\color{gray}חמשת מרכיבי הדרישה ודוגמה\end{center}

אין צורך לכתוב את חמשת המרכיבים האלה כמסמך רשמי; אפילו מעבר מהיר בראש מספיק. המפתח הוא לפתח הרגל של “לחשוב לפני שפותחים את הפה”. תגלו שכשאתם יכולים לענות בבהירות על חמש השאלות האלה, תיאור הדרישות שלכם הופך מעצמו למובנה.

\booksection{ב. אמנות פירוק הדרישות בשכבות}

בפרויקטים גדולים יותר, תיאור דרישה אחד אינו מספיק. עליכם לפרק דרישה גדולה לתת-דרישות קטנות שאפשר לאמת כל אחת בנפרד. התהליך נקרא “פירוק דרישות”, ועקרון הליבה שלו: כל תת-דרישה אמורה להיות מסוגלת להסתיים, להיבדק ולהימסר בעצמאות מלאה.

ניקח לדוגמה “לבנות אפליקציה אישית לניהול חשבונות” — דרך פירוק הגיונית תהיה:

\begin{enumerate}
  \item ממש תיעוד בסיסי של הכנסות והוצאות (הוספה, צפייה ומחיקה של רשומות);
  \item הוסף קטגוריות (לתייג כל רשומה כאוכל, תחבורה וכדומה);
  \item ממש סיכום חודשי (סטטיסטיקת הכנסות והוצאות לפי חודש);
  \item הוסף ויזואליזציה של נתונים (הצגת מבנה ההוצאות בגרפים).
\end{enumerate}

כל צעד הוא תוספת פונקציונלית שאפשר להריץ ולאמת בעצמאות.

הכלל הזהב לפירוק דרישות: בסוף כל שלב הפרויקט צריך להיות "משהו שאפשר להשתמש בו", גם אם התכונות פשוטות. אל תשאיר לעולם את הפרויקט במצב "הכול חצי-גמור וכלום לא עובד".

לפירוק כזה שני יתרונות. ראשית, הוא מאפשר לכם לראות התקדמות תמידית — בכל צעד שמושלם האפליקציה מרוויחה פיצ'ר נוסף, ומשוב חיובי זה קריטי בפרויקטים ארוכים (הרי כולם זקוקים לתחושת הישג). שנית, הוא מצמצם כל משימת AI לגודל מספיק קטן ומוריד את הסיכון ל“סטיית הבנה” — ככל שהמשימה קטנה יותר, ההבנה של ה-AI מדויקת יותר ואיכות התוצר גבוהה יותר.

\booksection{ג. תבנית כתיבת PRD ותרגול}

בפרויקטים גדולים יותר (למשל אפליקציית פול-סטאק עם פרונט ובק-אנד ומספר מודולי פונקציונליות), חמשת מרכיבי הדרישה עלולים שלא להספיק. אז צריך לכתוב מסמך דרישות פרויקט (PRD) רשמי יותר. נקודה שנוטים להבין לא נכון: ה-PRD אינו נכתב לצורכי ה-AI — הוא נכתב לצורכי עצמכם, כדי לסדר את המחשבות, ורק אחר כך תוכלו לספק ל-AI הוראות מדויקות יותר.

PRD טוב לא צריך להיות ארוך, אבל עליו לכלול שישה חלקים: סקירת הפרויקט (משפט או שניים שמסבירים מה בונים), רשימת פיצ'רים (מה כל נקודת פונקציונליות), פתרון טכני (איזו מחסנית טכנולוגיות ואיזו ארכיטקטורה בגדול), מודל נתונים (מבני הנתונים המרכזיים), תיאור ממשק (איך נראים המסכים המרכזיים) ותוכנית פיתוח (בכמה שלבים להשלים).

\begin{lstlisting}
# PRD לאפליקציית ניהול הוצאות אישיות
## תיאור הפרויקט
אפליקציה קלת משקל לניהול הוצאות אישיות: רישום הכנסות והוצאות, קטגוריות, סיכום חודשי והצגת נתונים חזותית.
## רשימת יכולות
1. רישום פעולה: סכום, סוג (הכנסה/הוצאה), קטגוריה, הערה, תאריך
2. ניהול קטגוריות: אוכל, תחבורה, קניות, פנאי, אחר – מוגדרות מראש
3. סיכום חודשי: סך הכנסות, סך הוצאות ויתרה לכל חודש
4. הצגת הוצאות: תרשים עוגה עם חלקה של כל קטגוריה
5. ייצוא נתונים: ייצוא ל-CSV
## פתרון טכני
- Frontend: React + TypeScript + Tailwind CSS
- אחסון: localStorage (בלי בק)
- גרפים: Chart.js
- Build: Vite
## מודל נתונים
Record {
  id: string
  amount: number
  type: 'income' | 'expense'
  category: string
  note: string
  date: string  // YYYY-MM-DD
}
## תוכנית פיתוח
Step 1: רישום בסיסי (CRUD)
Step 2: ניהול קטגוריות
Step 3: סיכום חודשי
Step 4: הצגת נתונים
Step 5: ייצוא נתונים
\end{lstlisting}

שימו לב לסגנון הכתיבה של PRD כזה: רשימות תמציתיות ופורמט מובנה, במקום נאומים ארוכים. לכל נקודת פונקציונליות הגדרה ברורה, ומבני הנתונים מתוארים בטיפוסים פשוטים. PRD בסגנון כזה ברור לא רק לכם; כשתספקו אותו ל-AI, גם הוא יבין כל נקודת דרישה במדויק.

\booksection{ד. תרגול מלא: הדרך השלמה מרעיון לתוכנית}

אין טעם בדיבור תאורטי — בואו נעבור על תהליך הנדסת הדרישות המלא באמצעות דוגמה אחת שלמה. נניח שהרעיון שלכם הוא: “רוצה לבנות מעקב למידה שמתעד מה למדתי בכל יום”.

צעד ראשון: עוברים על חמשת מרכיבי הדרישה.

\begin{itemize}
  \item Who: שימוש אישי.
  \item What: כלי שמתעד את תוכן הלמידה היומי, עוקב אחר משך הלמידה ומאפשר לסקור את התקדמות הלמידה.
  \item Why: רוצה לפתח הרגל של למידה מתמשכת ולראות את ההצטברות שלי.
  \item How: קודם גרסת דפדפן, והנתונים נשמרים מקומית בדפדפן.
  \item Constraint: פיתוח מאפס, השלמת הפונקציונליות המרכזית בתוך שבוע.
\end{itemize}

צעד שני: פירוק הדרישות.

\begin{enumerate}
  \item Step 1: מימוש "תיעוד פריטי למידה" (תאריך, תוכן, משך);
  \item Step 2: הוספת "סטטיסטיקת משך למידה" (היום, מצטבר שבועי);
  \item Step 3: מימוש תצוגת "לוח שנה של למידה" (כמו גרף התרומות של GitHub);
  \item Step 4: הוספת "סקירת למידה" (הצגה אקראית של רשומות למידה מיום עבר, לעזור בחזרה).
\end{enumerate}

צעד שלישי: כותבים PRD קצר. לא צריך להתאמץ — הפירוק שלמעלה כבר מהווה את גוף ה-PRD. מארגנים אותו לקובץ .md ושמים בתיקיית השורש של הפרויקט, כמקור עיון משותף לכם ול-AI.

צעד רביעי: מתחילים לפתח את השלב הראשון. ההוראה ל-AI יכולה להיות:

\begin{lstlisting}
אני צריך את היכולת הראשונה של עוקב למידה: רישום מפגשי למידה.
המשתמש יכול להוסיף רשומה עם תאריך, תוכן הלימוד ומשך (דקות).
הנתונים נשמרים ב-localStorage.
פרונט ב-React + TypeScript.
תן לי את המימוש המלא.
\end{lstlisting}

שימו לב למאפייני ההוראה הזאת: היא מתארת רק את מטרת השלב הנוכחי (פונקציונליות התיעוד), בלי לערב פיצ'רים אחרים; היא מפרטת את האילוצים הטכניים (React, TypeScript, localStorage); והיא נותנת מספיק מידע כדי שה-AI ידע “מה לעשות” ו“איך”, תוך שהמימוש הקונקרטי נותר בידי ה-AI. זו הוראת תכנות בעזרת AI איכותית, שעברה עיבוד של הנדסת דרישות.

\begin{tcbpractice}{תרגול מעשי: כתבו כרטיס חמישה מרכיבים ל"רעיון שחיבבתם"}
בחרו פרויקט קטן שהבשיל אצלכם זמן רב, והקדישו 5 דקות למלא כרטיס: Who / What / Why / How / Constraint — משפט אחד לכל אחד. ואז פרקו אותו ל-3-5 שלבים שכל אחד מהם “שמיש בסופו”. תגלו: ברגע שהרעיון יורד על הכתב, הוא עובר מ“חלום בהקיץ” ל“רשימת מטלות”.

\end{tcbpractice}

\begin{tcbsummary}{סיכום הפרק}
\begin{itemize}
  \item חמשת מרכיבי הדרישה: Who / What / Why / How / Constraint — לחשוב לפני שפותחים את הפה
  \item כלל הזהב של פירוק דרישות: אחרי כל שלב הפרויקט הוא "משהו שמיש"
  \item ה-PRD נכתב לעצמכם: סקירה, רשימת פיצ'רים, פתרון טכני, מודל נתונים, ממשק, תוכנית
  \item הוראה טובה = לתאר רק את השלב הנוכחי + לפרט אילוצים + לומר "מה" ולא "איך"
  \item תבנית התרגול המלאה: חמישה מרכיבים פירוק PRD הוראה, בארבעה צעדים
\end{itemize}

\end{tcbsummary}
